\chapter{Discussões e Considerações Finais}
\label{cap:cap6}

A discussão sobre a viabilidade econômica da arquitetura proposta revela-se indispensável para validar sua aplicabilidade no contexto de escassez de recursos do semiárido brasileiro. Ao contrastar o protótipo desenvolvido (fundamentado em hardware acessível e processamento na borda) com as onerosas soluções proprietárias de mercado e o atual regime de inspeções manuais, identifica-se um potencial disruptivo na estrutura de custos de monitoramento.

A premissa defendida aqui não busca encerrar o debate financeiro, mas sim evidenciar que a dependência exclusiva de visitas presenciais impõe um custo operacional linear e cumulativo, precária para a vigilância contínua de infraestruturas remotas como a Barragem de Oiticica. Nesse sentido, a solução automatizada apresenta-se como uma alternativa estratégica que, ao transformar o custo variável de deslocamento em um investimento fixo e reduzido de aquisição, sugere um caminho promissor para a democratização do monitoramento de segurança, viabilizando a proteção de ativos críticos mesmo sob severas restrições orçamentárias.
